% !TeX program = pdflatex
\documentclass[11pt,a4paper]{article}
\usepackage{../../recursos/latex/estilo-notas}

\begin{document}
\cabecalhoaula{03}{Seleção de Modelos e Validação Cruzada}{AME §1.4--1.5; ISLP cap.~5}

%==============================================================================
\section*{Objetivos da aula}
%==============================================================================
Ao final desta aula o aluno deve ser capaz de:
\begin{itemize}[leftmargin=1.4cm]
  \item explicar por que o erro de treino é um estimador \emph{otimista} do risco;
  \item descrever \emph{data splitting} (treino/validação) e estimar o risco com a
        parte de validação;
  \item definir \emph{validação cruzada} LOOCV e $k$-dobras e relacioná-las;
  \item discutir o balanço viés--variância \emph{na própria estimativa do risco} e
        a escolha de $k$;
  \item usar a validação cruzada para \textbf{selecionar hiperparâmetros} sem
        vazamento;
  \item descrever o \emph{bootstrap} para quantificar incerteza de uma estimativa.
\end{itemize}

\begin{emsala}
Esta é a aula \emph{operacional} mais importante do curso: é a ferramenta que
usaremos para escolher \emph{todos} os \emph{tuning parameters} daqui em diante
($\lambda$ do Lasso, $k$ do KNN, profundidade da árvore, $C$ da SVM\dots).
Pergunta de abertura: \emph{``se eu ajustar um polinômio de grau 50 a 50 pontos,
qual o erro no treino? E isso quer dizer que é um bom modelo?''}
\end{emsala}

%==============================================================================
\section{O problema: estimar o risco}
%==============================================================================
Da Aula~01, o risco de uma função de predição $g$ é
$\risco(g)=\E[L(Y,g(\X))]$, esperança sobre uma observação \emph{nova}. Para
\emph{selecionar} um modelo (escolher entre candidatos $g\in\mathcal{G}$, ou
escolher um hiperparâmetro) precisamos comparar riscos --- mas $\risco(g)$ é
desconhecido porque depende de $P$. Logo, \textbf{tudo se resume a estimar bem o
risco}.

\subsection{Por que não usar o erro de treino?}
O candidato ingênuo é o \emph{erro quadrático médio no treino} (risco empírico):
\begin{equation}\label{eq:eqm}
  \mathrm{EQM}(g) = \frac1n \sum_{i=1}^n \big(Y_i - g(\X_i)\big)^2 .
\end{equation}
O problema: $g$ foi \emph{escolhida para ajustar exatamente esses pontos}. O
EQM\eqref{eq:eqm} é um estimador \textbf{otimista} (enviesado para baixo) de
$\risco(g)$, e esse viés --- o \emph{otimismo} --- cresce com a complexidade do
modelo.

\begin{emsala}
Cuidado com a palavra, e vale escrever as duas no quadro lado a lado: o viés desta
frase é o do \emph{EQM como estimador do risco},
$\E[\mathrm{EQM}(g)]-\risco(g)$, que \emph{cresce} com a complexidade. Não é o viés
do \emph{modelo} da Aula~01, $\E_\Dados[\rhat(\x_0)]-r(\x_0)$, que \emph{cai} com a
complexidade. Um é o erro do modelo, o outro é o erro da régua que mede o modelo ---
e o modelo mais flexível é o que tem viés pequeno \emph{e} régua mais mentirosa.
Alguém vai perguntar; melhor antecipar.
\end{emsala}

\begin{atencao}
Levado ao extremo, minimizar o erro de treino leva ao superajuste perfeito: um
modelo flexível o bastante \emph{interpola} os dados (erro de treino zero, em
aritmética exata --- ver a §3 da \texttt{Aula prática 03} para o que o ponto
flutuante faz com isso) e generaliza pessimamente. \textbf{Nunca} selecione
modelos nem reporte desempenho com o erro de treino.
\end{atencao}

\begin{observacao}
Penalização (AIC, BIC) é uma alternativa: estima-se
$\risco(g)\approx \mathrm{EQM}(g)+P(g)$, com $P(g)$ crescente na complexidade
(p.ex.\ $P(g)\propto p/\widehat\sigma^2$ no AIC). Funciona, mas exige um modelo
probabilístico e contagem de ``parâmetros'' --- nem sempre disponível (KNN,
florestas). A validação cruzada é \emph{agnóstica ao modelo} e por isso a
preferimos.
\end{observacao}

%==============================================================================
\section{Data splitting (conjunto de teste)}
%==============================================================================
A ideia mais simples: separe os dados \emph{aleatoriamente} em treinamento
(p.ex.\ 70\%) e teste (p.ex.\ 30\%):
\[
  \underbrace{(\X_1,Y_1),\dots,(\X_n,Y_n)}_{\text{treinamento}}\;,\;
  \underbrace{(\X_{n+1},Y_{n+1}),\dots,(\X_{n+m},Y_{n+m})}_{\text{teste}}.
\]
Daqui para a frente $n$ é o tamanho do \emph{treinamento} e $m$ o do teste.
Ajuste $g$ \emph{só} no treinamento e estime o risco \emph{só} no teste:
\begin{equation}\label{eq:datasplit}
  \riscohat(g) = \frac{1}{m}\sum_{i=n+1}^{n+m} \big(Y_i - g(\X_i)\big)^2 .
\end{equation}
Como o teste não foi usado para estimar $g$, \eqref{eq:datasplit} é um
estimador \emph{consistente} de $\risco(g)$ pela lei dos grandes números. A
proporção também não é arbitrária: treinar um modelo exige mais informação do que
avaliá-lo, e por isso a fatia maior fica com o treinamento.

\begin{emsala}
Enfatize o \emph{aleatoriamente}: se o banco vier ordenado por alguma variável
(p.ex.\ por data ou por classe), uma divisão pelos índices originais cria
conjuntos não representativos. Em \texttt{scikit-learn} isso é o
\texttt{train\_test\_split(\dots, shuffle=True)} (padrão).
\end{emsala}

\begin{atencao}[Limitações do data splitting]
\begin{enumerate}[leftmargin=1.5cm,nosep]
  \item \textbf{Desperdício de dados:} $g$ é treinada com apenas $n$ das $n+m$
        observações que você coletou; com poucos dados isso prejudica o ajuste.
  \item \textbf{Variabilidade:} a estimativa depende da divisão sorteada --- outra
        semente dá outro número. Pode ser instável.
\end{enumerate}
A validação cruzada resolve os dois usando \emph{todo o conjunto de treinamento}
--- e é importante ver o que ela \emph{não} substitui. A divisão treino/teste
continua sendo o primeiro passo, e o teste continua guardado até o fim. O que a
validação cruzada dispensa é um \emph{segundo} corte, o que separaria uma parte do
treino para escolher o modelo: esse papel passa a ser feito pelas dobras, girando
dentro do treino.
\end{atencao}

%==============================================================================
\section{Validação cruzada}
%==============================================================================
\subsection{Leave-one-out (LOOCV)}
Deixe \emph{uma} observação de fora por vez. Seja $g_{-i}$ o modelo ajustado sem a
$i$-ésima observação. O estimador é
\begin{equation}\label{eq:loocv}
  \riscohat_{\text{LOO}}(g) = \frac1n \sum_{i=1}^n \big(Y_i - g_{-i}(\X_i)\big)^2 .
\end{equation}
Cada $Y_i$ é predito por um modelo que \emph{não} o viu. LOOCV é
\emph{aproximadamente não-viesado} para o risco esperado e não tem aleatoriedade de
divisão. O custo aparente é alto ($n$ ajustes), mas:

\begin{proposicao}[Atalho do LOOCV para modelos lineares]\label{prop:loo}
Para um ajuste linear (MQO, Ridge, \emph{splines}) com matriz de projeção
$\bm{H}=\bm{X}(\bm{X}^\top\bm{X})^{-1}\bm{X}^\top$, vale
\[
  \riscohat_{\text{LOO}}
   = \frac1n \sum_{i=1}^n \left( \frac{Y_i-\widehat{g}(\X_i)}{1-h_{ii}} \right)^2,
\]
onde $h_{ii}$ é o $i$-ésimo elemento diagonal de $\bm{H}$ (\emph{leverage}) e
$\widehat g$ é ajustado com \emph{todos} os dados.
\end{proposicao}
Ou seja, um \emph{único} ajuste basta. Esta é uma das fórmulas mais elegantes da
área: o LOOCV ``de graça''.

\subsection{Validação cruzada com $k$ dobras}
Divida \emph{o conjunto de treinamento} aleatoriamente em $k$ lotes (dobras)
disjuntos de tamanho $\approx n/k$, com índices $L_1,\dots,L_k$. O conjunto de teste
saiu na divisão anterior e não entra em dobra nenhuma. Para cada $j$, treine
$g_{-j}$ \emph{sem} o lote $L_j$ e avalie nesse lote:
\begin{equation}\label{eq:kfold}
  \riscohat_{k\text{-CV}}(g)
   = \frac1n \sum_{j=1}^{k} \sum_{i\in L_j} \big(Y_i - g_{-j}(\X_i)\big)^2 .
\end{equation}
Cada observação é usada para validação \emph{exatamente uma vez} e para treino
$k-1$ vezes. Quando $k=n$ recuperamos o LOOCV \eqref{eq:loocv}. Valores típicos:
$k=5$ ou $k=10$.

\begin{emsala}
Alguém que tenha lido o [ISLP] vai estranhar: a fórmula (5.3) de lá é
$\mathrm{CV}_{(k)}=\frac1k\sum_j \mathrm{EQM}_j$, a média dos EQM das dobras. É a
mesma coisa --- agrupando \eqref{eq:kfold} por dobra sai
$\sum_j (\abs{L_j}/n)\,\mathrm{EQM}_j$, a mesma média \emph{ponderada pelo tamanho
das dobras}. Coincidem exatamente quando $k$ divide $n$. A nossa pesa cada
observação igualmente; a do [ISLP], cada dobra --- e é a do [ISLP] que o
\texttt{cross\_val\_score} calcula. Quando $k$ não divide $n$, a §6 da
\texttt{Aula prática 03} mede a diferença: $0{,}9\%$ com $n=50$, $k=7$, e o grau
escolhido muda em 9 de 300 amostras. Daí a regra prática: com $n$ pequeno, escolha
um $k$ que divida $n$.
\end{emsala}

\begin{ideia}
A validação cruzada estima o risco \emph{de um procedimento} (``ajustar este tipo
de modelo a dados como estes''), não de um $g$ específico. Por isso, depois de
escolher o hiperparâmetro por CV, \textbf{reajustamos o modelo final usando todos
os dados}.
\end{ideia}

\subsection{Viés, variância e a escolha de $k$}
O \emph{argumento clássico} para escolher $k$ é ele mesmo um balanço
viés--variância, mas \emph{da estimativa do risco}:
\begin{itemize}[leftmargin=1.4cm]
  \item \textbf{$k$ grande (LOOCV):} cada $g_{-j}$ treina com $\approx n$ dados
        $\Rightarrow$ \emph{pouco viés}. Mas os $n$ modelos são quase idênticos e
        suas predições muito correlacionadas --- e a média de quantidades
        correlacionadas não se beneficia da lei dos grandes números como a de
        quantidades independentes. O argumento conclui: \emph{variância alta}.
        O custo também é alto (salvo o atalho da Prop.~\ref{prop:loo}).
  \item \textbf{$k$ pequeno (p.ex.\ $5$):} cada $g_{-j}$ treina com menos dados
        $\Rightarrow$ \emph{mais viés}, no sentido \emph{pessimista}: o procedimento
        avaliado não é bem o que você vai usar no final. Em troca, as dobras são
        mais independentes.
\end{itemize}

\begin{emsala}
Pergunte à turma o que acontece com o viés quando $k$ cresce, antes de responder.
O caminho é: cada dobra treina com $n-n/k$ observações, então $k$ maior aproxima o
procedimento avaliado daquele que será usado no fim --- e o viés, que é pessimista,
encolhe. A simulação da versão do aluno mede: $0{,}43$ em $k=2$ contra $0{,}006$ em
$k=10$.
\end{emsala}

\noindent Daí a conclusão prática: de $k=5$ em diante não há praticamente nada a
ganhar, e o custo continua subindo linearmente. Por isso $k=5$ ou $k=10$.

\begin{atencao}[A regra de ouro contra vazamento]
\textbf{Toda} etapa que ``aprende'' algo dos dados --- padronização, seleção de
variáveis, escolha de $\lambda$ --- deve ocorrer \emph{dentro} de cada dobra,
usando apenas o treino daquela dobra. Padronizar com a média/desvio de todo o
conjunto antes da CV é \emph{data leakage} e produz estimativas otimistas. (É por
isso que usamos \emph{pipelines}, Aula~07.)
\end{atencao}

\subsection{Selecionando hiperparâmetros}
O uso central: para cada valor de um hiperparâmetro $\theta$ (p.ex.\ $\lambda$ do
Lasso), calcule $\riscohat_{k\text{-CV}}(\theta)$ e escolha
$\widehat\theta=\argmin_\theta \riscohat_{k\text{-CV}}(\theta)$.

\begin{ideia}[o protocolo, em quatro passos]
Junte as peças e o procedimento inteiro fica assim:
\begin{enumerate}[leftmargin=1.6cm,nosep]
  \item \textbf{separe} o conjunto de teste (p.ex.\ 30\%) e guarde-o;
  \item \textbf{escolha} um ou mais bons candidatos por validação cruzada
        \emph{dentro do treino};
  \item \textbf{reajuste} o(s) candidato(s) em \emph{todo} o conjunto de
        treinamento --- as dobras serviram para escolher, não para entregar o
        modelo;
  \item \textbf{meça} o desempenho no conjunto de teste.
\end{enumerate}
O passo 3 é o que quase todo mundo esquece, e é ele que faz a CV valer a pena: o
modelo entregue usa \emph{todas} as observações de treino, não $k-1$ delas.
\end{ideia}

\begin{atencao}[Regra ``1-EP'']
Como $\riscohat_{k\text{-CV}}$ tem erro padrão, muitas vezes prefere-se a
\emph{regra de um erro-padrão}: escolher o modelo \emph{mais simples} cujo risco
estimado esteja a no máximo um erro-padrão do mínimo. Favorece parcimônia
(é o \texttt{lambda.1se} do \texttt{glmnet}).
\end{atencao}

\begin{atencao}[Validação \emph{aninhada}]
Se você usa CV para \emph{escolher} $\theta$ \emph{e} para \emph{reportar} o
desempenho, precisa de \textbf{CV aninhada}: um laço externo estima o desempenho,
um laço interno escolhe $\theta$. Usar a mesma CV para as duas coisas subestima o
risco.
\end{atencao}

%==============================================================================
\section{O bootstrap}
%==============================================================================
A validação cruzada estima o \emph{risco}; o \textbf{bootstrap} estima a
\emph{incerteza} (variância, viés, intervalos) de praticamente qualquer
estatística $\widehat\theta$ calculada a partir dos dados.

\begin{ideia}
Não conhecemos $P$, mas temos a amostra. O bootstrap usa a \emph{distribuição
empírica} (reamostrar dos dados \emph{com reposição}) como substituta de $P$,
imitando o ato de coletar novas amostras.
\end{ideia}

\noindent\textbf{Algoritmo.} Para $b=1,\dots,B$:
\begin{enumerate}[leftmargin=1.6cm,nosep]
  \item sorteie $n$ observações \emph{com reposição} da amostra original (amostra
        bootstrap $\Dados^{*b}$);
  \item calcule a estatística $\widehat\theta^{*b}$ em $\Dados^{*b}$.
\end{enumerate}
A dispersão de $\{\widehat\theta^{*1},\dots,\widehat\theta^{*B}\}$ estima a
distribuição amostral de $\widehat\theta$; por exemplo,
$\widehat{\mathrm{EP}}(\widehat\theta)$ é o desvio-padrão dessas réplicas, e um
intervalo de confiança de 95\% pode ser dado pelos percentis $2{,}5\%$ e $97{,}5\%$.

\begin{observacao}
Cada amostra bootstrap deixa de fora, em média, $\approx 1/e \approx 37\%$ das
observações originais (as \emph{out-of-bag}). Essa ideia reaparece no
\emph{bagging} e nas florestas aleatórias (Aula~06), onde as observações OOB dão
uma estimativa de risco ``de graça''.
\end{observacao}

%==============================================================================
\section{Prática em Python}
%==============================================================================
\begin{lstlisting}
from sklearn.model_selection import (train_test_split, cross_val_score, KFold,
                                     GridSearchCV)
from sklearn.pipeline import make_pipeline
from sklearn.preprocessing import StandardScaler
from sklearn.linear_model import Ridge
from sklearn.metrics import mean_squared_error

# 0. O PRIMEIRO passo e separar o teste. Dobra nenhuma vai toca-lo.
X_tr, X_te, y_tr, y_te = train_test_split(X, y, test_size=0.3, random_state=0)

modelo = make_pipeline(StandardScaler(), Ridge(alpha=1.0))
dobras = KFold(5, shuffle=True, random_state=0)

# 1. CV DENTRO DO TREINO, para estimar o risco de UM procedimento
#    (MSE -> negativo no sklearn)
scores = cross_val_score(modelo, X_tr, y_tr, cv=dobras,
                         scoring="neg_mean_squared_error")
print("risco estimado:", -scores.mean(), "+/-", scores.std())

# 2. Selecionar o hiperparametro lambda (=alpha) por CV, SEM vazamento
#    (o StandardScaler e reajustado dentro de cada dobra)
grade = {"ridge__alpha": [0.01, 0.1, 1, 10, 100]}
busca = GridSearchCV(modelo, grade, cv=dobras, scoring="neg_mean_squared_error")
busca.fit(X_tr, y_tr)                 # o teste continua guardado
print("melhor alpha:", busca.best_params_)

# 3. O vencedor ja saiu reajustado em TODO o treino (refit=True, o padrao).
#    So agora o teste e tocado, uma unica vez.
print("EQM no teste:", mean_squared_error(y_te, busca.predict(X_te)))
\end{lstlisting}

O notebook \texttt{Aula prática 03} implementa as $k$ dobras à mão, confere o atalho
do LOOCV pela alavancagem, mede o efeito de rodar \texttt{KFold} sem
\texttt{shuffle}, roda o protocolo de quatro passos contra a resposta conhecida e
fecha escolhendo o $\lambda$ de uma Ridge no \texttt{superconductivity.csv}.

%==============================================================================
\section*{Resumo}
%==============================================================================
\begin{itemize}[leftmargin=1.4cm]
  \item O erro de treino é otimista; estimar o risco exige dados \emph{não}
        usados no ajuste.
  \item \emph{Data splitting}: simples, mas desperdiça dados e é instável.
  \item \emph{LOOCV} \eqref{eq:loocv}: quase não-viesado; atalho fechado para
        modelos lineares (Prop.~\ref{prop:loo}).
  \item \emph{$k$-dobras} \eqref{eq:kfold}: usa toda a amostra; $k=5$ ou $10$
        equilibram viés/variância/custo.
  \item Use CV para escolher \emph{tuning parameters}; cuidado com \emph{vazamento}
        (faça tudo dentro da dobra) e com CV aninhada para reportar desempenho.
  \item \emph{Bootstrap}: reamostragem com reposição para medir incerteza de
        estimativas; conecta com OOB e \emph{bagging}.
\end{itemize}

\paragraph{Leitura recomendada.} [AME] §1.4 (função de risco e Teorema~1), §1.5.1
(\emph{data splitting} e validação cruzada); [ISLP] Capítulo~5 (§5.1 validação
cruzada, §5.2 \emph{bootstrap}).

\end{document}
